% sec:ad-forward — Forward-Mode AD and Dual Numbers
\section{Forward-Mode AD and Dual Numbers}
\label{sec:ad-forward}

Every evaluation of the geodesic right-hand side
(\cref{eq:geodesic-rhs}) requires the 40 independent
derivatives $\pd{\mu} g_{\alpha\beta}$ of the metric
tensor --- ten independent components differentiated with
respect to each of four coordinates.
Hand-coding these derivatives for each spacetime is tedious
and error-prone: the Kerr--Schild metric involves products of
transcendental functions and rational expressions in
Cartesian coordinates, and a single sign error in one of the
40 expressions produces trajectories that violate the
Hamiltonian constraint in ways that can be difficult to
diagnose.  Forward-mode automatic differentiation eliminates
the problem entirely:
\coderef{Spacetime}{AutoDiff}
by redefining arithmetic so that every
operation simultaneously computes its own derivative, the
integrator obtains exact derivatives as a byproduct of
evaluating the metric function, at a cost comparable to a
handful of ordinary evaluations.

\paragraph{Dual numbers.}

The idea is to extend the real numbers by an abstract
infinitesimal.  A \emph{dual number} is a pair
$a + b\,\varepsilon$, where $a$ and $b$ are real and the
symbol $\varepsilon$ satisfies
%
\begin{equation}
  \varepsilon^2 = 0 \,.
  \label{eq:dual-epsilon}
\end{equation}
%
This single algebraic rule --- the square of $\varepsilon$
vanishes, but $\varepsilon$ itself does not --- is the
engine of automatic differentiation.  The real part $a$
carries the function value; the dual part $b$ carries the
derivative.  The resulting algebraic structure
$\mathbb{D} = \{a + b\,\varepsilon \mid a, b \in \mathbb{R}\}$
forms a commutative ring under the arithmetic operations
defined below.
\histnote{Clifford introduced dual numbers in 1873, well
before their rediscovery for automatic differentiation in
the 1960s \cite{Griewank2008}.}

\paragraph{Arithmetic.}

The four arithmetic operations on dual numbers follow
directly from $\varepsilon^2 = 0$.  Addition is
component-wise:
%
\begin{equation}
  (a + b\,\varepsilon) + (c + d\,\varepsilon)
    = (a + c) + (b + d)\,\varepsilon \,.
  \label{eq:dual-add}
\end{equation}
%
The dual part $b + d$ is the sum of the two derivatives ---
the sum rule $(f + g)' = f' + g'$.
Multiplication expands as usual, with the $\varepsilon^2$
term vanishing:
%
\begin{equation}
  (a + b\,\varepsilon)\,(c + d\,\varepsilon)
    = ac + (ad + bc)\,\varepsilon \,.
  \label{eq:dual-mul}
\end{equation}
%
The dual part $ad + bc$ is precisely the product rule:
if $a = f(x)$, $b = f'(x)$, $c = g(x)$, $d = g'(x)$, then
$ad + bc = f(x)\,g'(x) + f'(x)\,g(x) = (fg)'(x)$.
Division follows from rationalising the denominator
(multiplying by $c - d\,\varepsilon$), provided $c \neq 0$:
%
\begin{equation}
  \frac{a + b\,\varepsilon}{c + d\,\varepsilon}
    = \frac{a}{c} + \frac{bc - ad}{c^2}\,\varepsilon \,,
  \label{eq:dual-div}
\end{equation}
%
and the dual part $(bc - ad)/c^2$ reproduces the quotient
rule.  Every arithmetic operation on dual numbers
automatically propagates the correct differentiation rule
through its dual part.
\physnote{Dual numbers resemble complex numbers ($i^2 = -1$),
but the nilpotent condition $\varepsilon^2 = 0$ makes them
fundamentally different: $\mathbb{D}$ contains zero divisors
(e.g.\ $\varepsilon \cdot \varepsilon = 0$) and is not a
field.  The nilpotency is exactly what truncates the Taylor
series to first order.}

\paragraph{Lifting elementary functions.}

Arithmetic covers polynomials, but the metric function
involves square roots, divisions, and --- for rotating
spacetimes --- trigonometric functions.  Each elementary
function is lifted to dual numbers by defining its action on
a dual argument as
%
\begin{equation}
  f(a + b\,\varepsilon)
    = f(a) + b\,f'(a)\,\varepsilon \,.
  \label{eq:dual-lift}
\end{equation}
%
This is the first-order Taylor expansion, but here it is
\emph{exact} rather than approximate: the nilpotency of
$\varepsilon$ kills every higher-order term, so there is no
remainder.  For example,
$\sin(a + b\,\varepsilon) = \sin a + b\cos a\;\varepsilon$,
and
$\sqrt{a + b\,\varepsilon}
  = \sqrt{a} + \frac{b}{2\sqrt{a}}\,\varepsilon$.
Each elementary function needs its own lifting rule, but the
rule is mechanical: pair the standard value with the standard
derivative.  An AD library provides these liftings for every
function in the language's numeric type class.
\warnnote{The lifting \cref{eq:dual-lift} requires $f$ to
be differentiable at $a$.  Functions with branch points or
discontinuities --- such as $\abs{x}$ at $x = 0$ --- produce
undefined or incorrect dual parts at those points.  The
metric functions used by the ray tracer are smooth everywhere
except at the physical singularity, which the integrator
never reaches.}

\paragraph{The chain rule for free.}

The power of dual-number arithmetic is that the chain rule
emerges from composition without any explicit
differentiation.  Seeding with unit dual part ($b = 1$)
and applying $g$ after $f$ gives
%
\begin{equation}
  g\bigl(f(a + \varepsilon)\bigr)
    = g\bigl(f(a) + f'(a)\,\varepsilon\bigr)
    = g(f(a)) + f'(a)\,g'(f(a))\,\varepsilon \,.
  \label{eq:dual-chain}
\end{equation}
%
The dual part of the result is
$f'(a)\,g'(f(a)) = (g \circ f)'(a)$ --- exactly the chain
rule.  No matter how deeply nested the composition, each
intermediate result carries its derivative forward, and the
chain rule accumulates automatically through the sequence of
dual-number operations.  A metric function that calls
\texttt{sqrt}, multiplies by a coordinate, divides by a sum,
and adds a constant produces the exact derivative of the
full expression as a byproduct of evaluation.

\paragraph{Seeding and directional derivatives.}

Given a function $f : \mathbb{R}^n \to \mathbb{R}^m$,
dual numbers compute directional derivatives.  To
obtain $\partial f / \partial x^i$, \emph{seed} the $i$-th
input with dual part one and all other inputs with dual part
zero:
%
\begin{equation}
  \tilde{x}^j = x^j + \delta^j{}_i\,\varepsilon \,.
  \label{eq:dual-seed}
\end{equation}
%
Evaluating $f(\tilde{x}^0, \ldots, \tilde{x}^{n-1})$ then
yields all $m$ output components with their partial
derivatives with respect to $x^i$ stored in the dual parts.
For the metric function --- which maps four coordinates to
ten independent metric components --- a single seeded
evaluation produces one column of the $10 \times 4$ Jacobian
matrix $\partial g_{\alpha\beta} / \partial x^\mu$.  Four
evaluations, one per coordinate direction, produce the
complete set of 40 derivatives.

\paragraph{The codebase implementation.}

The event-horizon codebase avoids explicit dual-number
types in user code through parametric polymorphism.  Every
metric function has a polymorphic type
signature --- it accepts any numeric type that supports the
standard floating-point operations, not just
\texttt{Double}.  When called with ordinary double-precision
values, the function evaluates the metric normally.  When
called with the AD library's dual-number type, the same
code --- unchanged, not rewritten --- propagates derivatives
alongside values.
\implnote{The rank-2 polymorphism (\texttt{forall a.}) ensures
that the caller, not the metric function, chooses the numeric
type.  The \texttt{ad} library's \texttt{jacobian} function
exploits this by internally substituting a dual-number type.}

This pattern is wrapped into a single call by the function
\texttt{metricDerivativesAD} in
\codemodule{Spacetime.AutoDiff}:
%
\begin{lstlisting}[language=Haskell]
metricDerivativesAD
  :: (forall a. Floating a => [a] -> [a])
  -> [Double]
  -> (Sym4 'Covariant, Sym4 'Covariant,
      Sym4 'Covariant, Sym4 'Covariant)
\end{lstlisting}
%
It accepts a polymorphic metric function and a coordinate
point, and returns four \texttt{Sym4} tensors --- one per
coordinate direction --- containing $\pd{\mu} g_{\alpha\beta}$
for $\mu = 0, 1, 2, 3$.  Internally, it calls the
\texttt{ad} library's \texttt{jacobian} function, which
performs the four seeded evaluations of
\cref{eq:dual-seed} and collects the results into a
$10 \times 4$ matrix, from which each column is packed into a
symmetric tensor.

Physical parameters such as the black hole mass $M$ or spin
$a$ enter through partial application before the metric
function reaches the AD machinery.  The Schwarzschild metric
function, for example, has type
\texttt{Floating a => a -> [a] -> [a]}, with mass as its
first argument.  The Christoffel computation partially
applies the mass via \texttt{realToFrac}~$M$,
producing a closure of type
\texttt{forall a. Floating a => [a] -> [a]} that treats
$M$ as a constant during differentiation --- exactly the
correct mathematical behaviour, since the derivatives
$\pd{\mu} g_{\alpha\beta}$ are taken with respect to
coordinates, not parameters.

\paragraph{Cost.}

All 40 metric derivatives cost roughly eight times a single
metric evaluation --- a factor that sounds expensive but is
modest in practice.  The factor arises because each of the
four seeded passes costs
approximately twice a standard evaluation (every
floating-point operation is paired with a derivative
operation), giving $4 \times 2 = 8$.  The Dormand--Prince
integrator (\cref{sec:rk-dormand-prince}) calls the
right-hand side six times per accepted step (after FSAL
reuse), so the derivative computation adds roughly $48$
metric evaluations per step --- far less than the cost of
implementing, debugging, and maintaining hand-coded
derivative expressions for each spacetime.
\xrefnote{A quantitative comparison of AD, finite-difference,
and analytic derivatives --- including accuracy and wall-clock
timing --- appears in \cref{sec:ad-comparison}.}

The accuracy payoff is equally significant: AD produces
derivatives correct to machine precision, with no truncation
error from finite-difference step sizes and no risk of
transcription errors in analytic expressions.  For a
reference ray tracer where correctness matters more than raw
speed, this tradeoff is decisive.
